% GENERATED FILE. Edit canonical lessons, then run npm run generate:books.
% canonical-source-hash: fnv1a64:183c63db05592043
% canonical-lessons: SA-W258-a2-connected-composition

\chapter{A connected A2 message}
\label{ch:sa-a-connected-a2-message}
\hlchaptermodality{258}

\begin{chapteropening}
\textbf{By the end of this chapter:} \emph{I can plan and write a connected 70–90 word Sanskrit message in Devanagari for a named reader and practical purpose.}

Complete chapter 258 by writing a model-free 70–90 word connected Sanskrit message for a named reader and practical purpose.
\end{chapteropening}

\section[(connected A2 writing)]{A connected message about your familiar day}

\label{lesson:SA-W258-a2-connected-composition}



\begin{quote}
A connected message is not a pile of vocabulary and not one long sentence. It is a short route for one reader. Before writing, draw eight empty boxes and give them these jobs:

\begin{enumerate}
\raggedright
  \item reader and greeting;
  \item reason for writing;
  \item where or when your day begins;
  \item two familiar actions;
  \item what happens next and a second time anchor;
  \item one preference and one reason;
  \item one direct question about tomorrow;
  \item a close.
\end{enumerate}

Write only one or two English cue words in each box. Do not draft Sanskrit yet.
\end{quote}

\subsection*{Writing — connected composition}
There is no model, Sanskrit sentence bank, or copyable answer on this page.

\textbf{Reader:} your friend Kiran. \textbf{Purpose:} help Kiran understand your familiar day and choose a good time to meet tomorrow.

Write \textbf{70–90 words in Sanskrit using Devanagari}. Use every planning box. Include at least four familiar actions, two time anchors, one sequence link, one preference, one reason, and one direct question for Kiran. Make every detail help with tomorrow's meeting. The instructions give content, not sentences: choose only words and forms this book has taught, and connect them yourself. Use only external-sandhi changes that the book has taught; simplify an uncertain boundary instead of inventing a join.

This first connected attempt is untimed. Put away romanization, dictionaries, grammar references, translators, generative assistants, and spell-checkers. Do not copy from an earlier lesson. Ordinary legible Devanagari handwriting is enough.

When the whole message exists, count its words once. If it is short, add one useful taught detail to a planning box. If it is long, remove a detail that does not help Kiran. Do not replace the message with a memorised model.

\begin{tcolorbox}[breakable,colback=teal!4,colframe=teal!35!black,title={Before you move on}]
Review the same attempt in four separate passes:

\begin{enumerate}
\raggedright
  \item \textbf{Task fulfilment and relevant content} — Kiran, the familiar day, tomorrow's meeting, and every requested content point are present.
  \item \textbf{Comprehensibility and organisation} — the eight jobs form a route rather than a list of unrelated words or lines.
  \item \textbf{Sanskrit vocabulary and grammar} — every form and external-sandhi change comes from the taught Sanskrit track; a simpler familiar boundary is better than an invented join.
  \item \textbf{Devanagari orthographic control} — word-final visarga or anusvara, continuous śirorekhā, word boundaries, and danda or other punctuation remain readable.
\end{enumerate}

Mark one strength and one next step. Repair only the next step; keep the original attempt so the next 70–90-word message can show real growth.

This is one untimed curriculum practice attempt, not the complete 30-minute A2 writing paper and not evidence that either timed mock has been passed. Closing the machine-measured connected-composition stage does not prove A2 attainment or readiness, learner proof, independent review, standard setting, or complete external-sandhi mastery.
\end{tcolorbox}
